\documentclass[12pt,a4paper]{article}
\usepackage[utf8]{vietnam}
\usepackage{amsmath,amssymb}
\usepackage{tabularx}
\usepackage{tikz}
\usetikzlibrary{calc,angles,quotes,arrows.meta,patterns}
\usepackage{xcolor}
\usepackage{enumitem}
\usepackage{fancyhdr}
\usepackage{graphicx}
\usepackage[top=0.9cm, bottom=0.9cm, left=1.2cm, right=1.2cm]{geometry}
\usepackage[loigiai]{ex_test}
\SAOPTN{widthSA=3.2,heightSA=0.62}

\setlength{\headheight}{15pt}
\setlength{\parskip}{0pt}

\pagestyle{fancy}
\fancyhf{}
\rhead{\textit{Môn Toán 8 -- Hướng dẫn giải Đề số 03}}
\lhead{\textit{Bộ đề ôn tập Giữa học kì I}}
\cfoot{\thepage}

\begin{document}

\begin{center}
  \textbf{\Large HƯỚNG DẪN GIẢI CHI TIẾT VÀ BẢNG ĐÁP ÁN}\\
  \textbf{\large ĐỀ ÔN TẬP GIỮA HỌC KỲ I -- ĐỀ SỐ 03}\\
  \textit{Môn: Toán 8 -- Năm học 2026 -- 2027}
\end{center}
\vspace{2mm}

% --- BẢNG ĐÁP ÁN TRẮC NGHIỆM TÓM TẮT ---
\noindent\textbf{\large A. BẢNG ĐÁP ÁN TRẮC NGHIỆM TỔNG HỢP}
\vspace{1mm}

\noindent\textbf{1. Phần I (Mỗi câu trả lời đúng được 0,25 điểm):}
\begin{center}
\begin{tabular}{|c|c|c|c|c|c|c|c|c|c|c|c|c|}
\hline
\textbf{Câu} & \textbf{1} & \textbf{2} & \textbf{3} & \textbf{4} & \textbf{5} & \textbf{6} & \textbf{7} & \textbf{8} & \textbf{9} & \textbf{10} & \textbf{11} & \textbf{12} \\
\hline
\textbf{Đ/A} & \textbf{B} & \textbf{C} & \textbf{B} & \textbf{D} & \textbf{A} & \textbf{D} & \textbf{D} & \textbf{A} & \textbf{C} & \textbf{D} & \textbf{A} & \textbf{B} \\
\hline
\end{tabular}
\end{center}
\vspace{0.5mm}

\noindent\textbf{2. Phần II (Đúng 1 ý: 0,1đ; 2 ý: 0,25đ; 3 ý: 0,5đ; 4 ý: 1,0đ):}
\begin{center}
\begin{tabular}{|c|c|c|c|c||c|c|c|c|c|}
\hline
\textbf{Câu} & \textbf{a)} & \textbf{b)} & \textbf{c)} & \textbf{d)} & \textbf{Câu} & \textbf{a)} & \textbf{b)} & \textbf{c)} & \textbf{d)} \\
\hline
\textbf{Câu 1} & \textbf{Đ} & \textbf{Đ} & \textbf{Đ} & \textbf{S} & \textbf{Câu 2} & \textbf{Đ} & \textbf{Đ} & \textbf{Đ} & \textbf{S} \\
\hline
\end{tabular}
\end{center}
\vspace{0.5mm}

\noindent\textbf{3. Phần III (Mỗi câu trả lời đúng được 1,0 điểm):}
\begin{center}
\begin{tabular}{|c|c|c|}
\hline
\textbf{Câu} & \textbf{Câu 1} & \textbf{Câu 2} \\
\hline
\textbf{Đáp án} & $\mathbf{x - 3}$ & $\mathbf{(x - y + 1)(x + y + 1)}$ \\
\hline
\end{tabular}
\end{center}
\vspace{1.5mm}

% --- B. LỜI GIẢI CHI TIẾT ---
\noindent\textbf{\large B. LỜI GIẢI CHI TIẾT CÁC CÂU HỎI}
\vspace{1mm}
\begin{ex}
Biểu thức nào sau đây là đơn thức?
\choice
{$x^3(1 + x)$}
{\True $3xy^2z^5$}
{$(xy + z)t$}
{$\dfrac{2}{x} + y$}
\loigiai{
Đơn thức là biểu thức đại số chỉ gồm một số, hoặc một biến, hoặc một tích giữa các số và các biến. Biểu thức $3xy^2z^5$ là đơn thức.
}
\end{ex}

\begin{ex}
Cặp đơn thức nào sau đây là đồng dạng với nhau?
\choice
{$\dfrac{1}{2}x^2y$ và $2xy^2$}
{$3xy$ và $-5x$}
{\True $x^2y^4$ và $2x^2y^4$}
{$-3xy$ và $-3xy^2$}
\loigiai{
Hai đơn thức đồng dạng có hệ số khác 0 và có cùng phần biến. Cặp đơn thức $x^2y^4$ và $2x^2y^4$ có cùng phần biến $x^2y^4$.
}
\end{ex}

\begin{ex}
Thực hiện phép chia đơn thức $5x^2y^2z$ cho $5xyz$ ta được kết quả là:
\choice
{$5xy$}
{\True $xy$}
{$yz$}
{$25x^3y^3z^2$}
\loigiai{
Ta có: $5x^2y^2z : 5xyz = (5 : 5)(x^2 : x)(y^2 : y)(z : z) = xy$.
}
\end{ex}

\begin{ex}
Hằng đẳng thức $A^2 - B^2 = (A - B)(A + B)$ có tên là:
\choice
{Bình phương của một tổng}
{Bình phương của một hiệu}
{Tổng hai bình phương}
{\True Hiệu hai bình phương}
\loigiai{
Công thức $A^2 - B^2 = (A - B)(A + B)$ là hằng đẳng thức hiệu hai bình phương.
}
\end{ex}

\begin{ex}
Tứ giác nào dưới đây là hình thoi?
\choice
{\True Tứ giác có 4 cạnh bằng nhau}
{Tứ giác có hai đường chéo vuông góc}
{Hình bình hành có hai đường chéo bằng nhau}
{Tứ giác có hai đường chéo cắt nhau tại trung điểm mỗi đường}
\loigiai{
Theo định nghĩa, hình thoi là tứ giác có 4 cạnh bằng nhau.
}
\end{ex}

\begin{ex}
Hai đường chéo của hình chữ nhật có tính chất nào sau đây?
\choice
{Chúng vuông góc với nhau}
{Chúng bằng nhau}
{Chúng cắt nhau tại trung điểm của mỗi đường}
{\True Chúng bằng nhau và cắt nhau tại trung điểm của mỗi đường}
\loigiai{
Trong hình chữ nhật, hai đường chéo bằng nhau và cắt nhau tại trung điểm của mỗi đường.
}
\end{ex}

\begin{ex}
Hình bình hành $ABCD$ có $\widehat{A} = 65^\circ$. Số đo góc $\widehat{C}$ là:
\choice
{$125^\circ$}
{$95^\circ$}
{$115^\circ$}
{\True $65^\circ$}
\loigiai{
Trong hình bình hành, hai góc đối diện bằng nhau nên $\widehat{C} = \widehat{A} = 65^\circ$.
}
\end{ex}

\begin{ex}
Cho tam giác $ABC$ vuông tại $A$ có $M$ là trung điểm của $BC$, biết $BC = 10\text{ cm}$. Khi đó độ dài đoạn thẳng $AM$ bằng:
\choice
{\True $5\text{ cm}$}
{$10\text{ cm}$}
{$15\text{ cm}$}
{$20\text{ cm}$}
\loigiai{
Trong tam giác vuông, đường trung tuyến ứng với cạnh huyền bằng một nửa cạnh huyền: $AM = \dfrac{1}{2}BC = \dfrac{1}{2} \cdot 10 = 5\text{ cm}$.
}
\end{ex}

\begin{ex}
Cho các tam giác có độ dài ba cạnh như sau. Tam giác nào là tam giác vuông?
\choice
{$15\text{ cm}; 8\text{ cm}; 10\text{ cm}$}
{$6\text{ dm}; 8\text{ dm}; 9\text{ dm}$}
{\True $5\text{ cm}; 4\text{ cm}; 3\text{ cm}$}
{$2\text{ cm}; 3\text{ cm}; 4\text{ cm}$}
\loigiai{
Theo định lý Pythagoras đảo: $3^2 + 4^2 = 9 + 16 = 25 = 5^2$, do đó tam giác có độ dài ba cạnh $3\text{ cm}, 4\text{ cm}, 5\text{ cm}$ là tam giác vuông.
}
\end{ex}

\begin{ex}
Cho hình chữ nhật $ABCD$. Đoạn thẳng $AC$ bằng đoạn thẳng nào sau đây?
\choice
{$AB$}
{$BC$}
{$AD$}
{\True $BD$}
\loigiai{
Trong hình chữ nhật $ABCD$, hai đường chéo $AC$ và $BD$ bằng nhau: $AC = BD$.
}
\end{ex}

\begin{ex}
Cho tứ giác $ABCD$, trong đó $\widehat{A} + \widehat{B} = 140^\circ$. Tổng $\widehat{C} + \widehat{D}$ bằng:
\choice
{\True $220^\circ$}
{$200^\circ$}
{$160^\circ$}
{$130^\circ$}
\loigiai{
Tổng 4 góc của tứ giác bằng $360^\circ$. Do đó: $\widehat{C} + \widehat{D} = 360^\circ - (\widehat{A} + \widehat{B}) = 360^\circ - 140^\circ = 220^\circ$.
}
\end{ex}

\begin{ex}
Thể tích của một hình chóp tam giác đều có diện tích đáy bằng $15\text{ cm}^2$ và chiều cao bằng $6\text{ cm}$ là:
\choice
{$90\text{ cm}^3$}
{\True $30\text{ cm}^3$}
{$45\text{ cm}^3$}
{$15\text{ cm}^3$}
\loigiai{
Thể tích hình chóp là: $V = \dfrac{1}{3} S_{\text{đáy}} \cdot h = \dfrac{1}{3} \cdot 15 \cdot 6 = 30\text{ (cm}^3\text{)}$.
}
\end{ex}


\setcounter{ex}{0}
\noindent\textbf{\large PHẦN II. CÂU TRẮC NGHIỆM ĐÚNG SAI}
\vspace{1mm}
\begin{ex}
Cho hai biểu thức $A = (2x - 1)(x + 3) - 2x(x + 2)$ và $B = (6x^3y^2 - 9x^2y^3) : 3x^2y$. Xét tính đúng sai của các mệnh đề sau:
\choiceTF
{\True Khai triển tích $(2x - 1)(x + 3)$ ta được $2x^2 + 5x - 3$}
{\True Kết quả rút gọn của biểu thức $A$ là $x - 3$}
{\True Kết quả của phép chia biểu thức $B$ là $2xy - 3y^2$}
{Tại $x = 2$ và $y = 1$, giá trị của biểu thức $A$ gấp 2 lần giá trị của biểu thức $B$}
\loigiai{
\textbf{a)} Đúng. $(2x - 1)(x + 3) = 2x^2 + 6x - x - 3 = 2x^2 + 5x - 3$.

\textbf{b)} Đúng. $A = (2x^2 + 5x - 3) - (2x^2 + 4x) = x - 3$.

\textbf{c)} Đúng. $B = (6x^3y^2 - 9x^2y^3) : 3x^2y = 2xy - 3y^2$.

\textbf{d)} Sai. Tại $x = 2, y = 1$: $A = 2 - 3 = -1$ và $B = 2 \cdot 2 \cdot 1 - 3 \cdot 1^2 = 1$. Vì $-1 \ne 2 \cdot 1$ nên khẳng định là sai.
}
\end{ex}

\begin{ex}
Cho hai đa thức $P = x^2 - 6x + 9 - y^2$ và $Q = 8x^3 - 27$. Xét tính đúng sai của các mệnh đề sau:
\choiceTF
{\True Biến đổi $x^2 - 6x + 9$ thành bình phương của một hiệu ta được $(x - 3)^2$}
{\True Phân tích đa thức $P$ thành nhân tử ta được $(x - y - 3)(x + y - 3)$}
{\True Phân tích đa thức $Q$ thành nhân tử theo hằng đẳng thức hiệu hai lập phương ta được $(2x - 3)(4x^2 + 6x + 9)$}
{Tại $x = 5$ và $y = 2$, giá trị của đa thức $P$ bằng 8}
\loigiai{
\textbf{a)} Đúng. $x^2 - 6x + 9 = (x - 3)^2$.

\textbf{b)} Đúng. $P = (x - 3)^2 - y^2 = (x - 3 - y)(x - 3 + y) = (x - y - 3)(x + y - 3)$.

\textbf{c)} Đúng. $Q = (2x)^3 - 3^3 = (2x - 3)(4x^2 + 6x + 9)$.

\textbf{d)} Sai. Thay $x = 5, y = 2$ vào $P$: $P = (5 - 3)^2 - 2^2 = 2^2 - 4 = 0 \ne 8$.
}
\end{ex}


\setcounter{ex}{0}
\noindent\textbf{\large PHẦN III. CÂU HỎI TRẢ LỜI NGẮN}
\vspace{1mm}
\begin{ex}
Rút gọn biểu thức: $A = (2x - 1)(x + 3) - 2x(x + 2)$. \shortans{$x - 3$}
\loigiai{
Ta có:
$$A = (2x^2 + 6x - x - 3) - (2x^2 + 4x) = 2x^2 + 5x - 3 - 2x^2 - 4x = x - 3.$$
Đáp số: $x - 3$.
}
\end{ex}

\begin{ex}
Phân tích đa thức sau thành nhân tử: $x^2 - y^2 + 2x + 1$. \shortans{$(x - y + 1)(x + y + 1)$}
\loigiai{
Nhóm các hạng tử tạo hằng đẳng thức:
$$x^2 + 2x + 1 - y^2 = (x + 1)^2 - y^2 = (x + 1 - y)(x + 1 + y) = (x - y + 1)(x + y + 1).$$
Đáp số: $(x - y + 1)(x + y + 1)$.
}
\end{ex}


\setcounter{ex}{0}
\noindent\textbf{\large PHẦN IV. TỰ LUẬN (HƯỚNG DẪN GIẢI VÀ THANG ĐIỂM)}
\vspace{1mm}
\begin{ex} (2,0 điểm)
Thực hiện phép tính:
\begin{enumerate}[label=\textbf{\alph*)}, leftmargin=0.8cm, itemsep=0pt, topsep=1pt]
\item (0,5 điểm) Thu gọn và tính giá trị biểu thức $A = 5x^2y - 2xy + 3x^2y - xy$ tại $x = -1, y = 2$.
\item (0,75 điểm) Thực hiện phép nhân: $B = 3xy(2x^2 - 5y^3)$.
\item (0,75 điểm) Thực hiện phép chia: $C = (6x^3y^2 - 9x^2y^3) : 3x^2y^2$.
\end{enumerate}
\loigiai{
\textbf{a)} \textit (0,5 điểm): Thu gọn đa thức: $A = (5x^2y + 3x^2y) + (-2xy - xy) = 8x^2y - 3xy$.
Thay $x = -1, y = 2$: $A = 8(-1)^2(2) - 3(-1)(2) = 16 + 6 = 22$.

\textbf{b)} \textit (0,75 điểm): $B = 3xy \cdot 2x^2 - 3xy \cdot 5y^3 = 6x^3y - 15xy^4$.

\textbf{c)} \textit (0,75 điểm): $C = 6x^3y^2 : 3x^2y^2 - 9x^2y^3 : 3x^2y^2 = 2x - 3y$.


}
\end{ex}
\vspace{0.5mm}

\begin{ex} (1,5 điểm)
Cho hai đa thức: $P = 2x^2y + 5x^3 - 2xy^2 + 6$ và $Q = x^3 + xy^2 - x^2y - 2$.
\begin{enumerate}[label=\textbf{\alph*)}, leftmargin=0.8cm, itemsep=0pt, topsep=1pt]
\item (0,75 điểm) Tính tổng của hai đa thức $P$ và $Q$.
\item (0,75 điểm) Tìm đa thức $N$ sao cho $Q = P + N$.
\end{enumerate}
\loigiai{
\textbf{a)} \textit (0,75 điểm): $\begin{aligned}[t] P + Q &= (5x^3 + x^3) + (2x^2y - x^2y) + (-2xy^2 + xy^2) + (6 - 2) \\ &= 6x^3 + x^2y - xy^2 + 4. \end{aligned}$

\textbf{b)} \textit (0,75 điểm): Ta có $N = Q - P$:
$$\begin{aligned} N &= (x^3 + xy^2 - x^2y - 2) - (2x^2y + 5x^3 - 2xy^2 + 6) \\ &= x^3 + xy^2 - x^2y - 2 - 2x^2y - 5x^3 + 2xy^2 - 6 \\ &= -4x^3 - 3x^2y + 3xy^2 - 8. \end{aligned}$$


}
\end{ex}
\vspace{0.5mm}

\begin{ex} (1,0 điểm)
Khu vườn của nhà bác Hoa có dạng hình vuông. Bác Hoa muốn dành một mảnh đất hình chữ nhật ở góc khu vườn làm nhà để dụng cụ làm vườn (hình vẽ).
\begin{center}
\begin{tikzpicture}[scale=0.8, baseline=(current bounding box.center)]
  \draw[thick, fill=green!10] (0,0) rectangle (4,4);
  \draw[thick, fill=orange!25] (0,0) rectangle (2.2,1.3);
  \node at (1.1,0.65) {\small Nhà kho};
  \node at (2.5,2.6) {\small Khu vườn};
  \draw[<->] (0,-0.3) -- (4,-0.3) node[midway, below] {\small $x$};
  \draw[<->] (-0.3,0) -- (-0.3,4) node[midway, left] {\small $x$};
  \draw[<->] (2.2,-0.3) -- (4,-0.3) node[midway, below] {\small $15\text{ m}$};
  \draw[<->] (-0.3,1.3) -- (-0.3,4) node[midway, left] {\small $25\text{ m}$};
\end{tikzpicture}
\end{center}
\begin{enumerate}[label=\textbf{\alph*)}, leftmargin=0.8cm, itemsep=0pt, topsep=1pt]
\item (0,5 điểm) Viết đa thức biểu thị chu vi của mảnh đất làm nhà.
\item (0,5 điểm) Biết chu vi của mảnh đất làm nhà bằng $40\text{ m}$. Tính diện tích khu vườn hình vuông ban đầu.
\end{enumerate}
\loigiai{
\textbf{a)} \textit (0,5 điểm): Mảnh đất hình chữ nhật có chiều dài là $x - 15\text{ (m)}$ và chiều rộng là $x - 25\text{ (m)}$.
Đa thức biểu thị chu vi mảnh đất là:
$$C = 2 \cdot [(x - 15) + (x - 25)] = 2(2x - 40) = 4x - 80\text{ (m)}.$$

\textbf{b)} \textit (0,5 điểm): Theo bài ra: $4x - 80 = 40 \Rightarrow 4x = 120 \Rightarrow x = 30\text{ (m)}$.
Cạnh khu vườn hình vuông ban đầu là $30\text{ m}$.
Diện tích khu vườn ban đầu là: $S = 30 \times 30 = 900\text{ (m}^2\text{)}$.


}
\end{ex}
\vspace{0.5mm}

\begin{ex} (0,5 điểm)
Cho các số thực $x, y$ thỏa mãn $x + y = 1$. Tìm giá trị nhỏ nhất của biểu thức $A = x^3 + y^3 + xy$.
\begin{enumerate}[label=\textbf{\alph*)}, leftmargin=0.8cm, itemsep=0pt, topsep=1pt]
\item (0,5 điểm) Tìm GTNN của $A$.
\end{enumerate}
\loigiai{
\textbf{a)} \textit (0,5 điểm): Biến đổi biểu thức $A$:
$$\begin{aligned} A &= (x + y)(x^2 - xy + y^2) + xy \\ &= 1 \cdot (x^2 - xy + y^2) + xy \\ &= x^2 + y^2. \end{aligned}$$
Thay $y = 1 - x$ vào biểu thức, ta được:
$$\begin{aligned} A &= x^2 + (1 - x)^2 = x^2 + 1 - 2x + x^2 \\ &= 2x^2 - 2x + 1 = 2\left(x^2 - x + \dfrac{1}{4}\right) + \dfrac{1}{2} \\ &= 2\left(x - \dfrac{1}{2}\right)^2 + \dfrac{1}{2}. \end{aligned}$$
Vì $2\left(x - \dfrac{1}{2}\right)^2 \ge 0$ với mọi $x$, nên $A \ge \dfrac{1}{2}$.
Dấu "$="$ xảy ra khi $x = \dfrac{1}{2} \Rightarrow y = 1 - \dfrac{1}{2} = \dfrac{1}{2}$.
Vậy giá trị nhỏ nhất của $A$ là $\dfrac{1}{2}$, đạt được khi $x = y = \dfrac{1}{2}$.


}
\end{ex}
\vspace{0.5mm}


\enlargethispage{1.5cm}
\begin{center}
\rule{4cm}{0.4pt}\\
\textit{--- HẾT ---}
\end{center}
\end{document}
